\documentclass[../readings.tex]{subfiles}

\begin{document}

\subsection{The Lanczos Algorithm}
\label{sub:lanczos-theory}

\textBox{%
The Lanczos algorithm is the symmetric specialization of the Arnoldi iteration for large, sparse real symmetric matrices $\bA = \bA^T \in \fR^{n \times n}$. When $\bA$ is symmetric, the upper Hessenberg projection matrix becomes symmetric tridiagonal, collapsing the $k$-term Gram-Schmidt orthogonalization into a three-term recurrence involving only the two immediately preceding basis vectors.
}

\textBox{%
By projecting the large $n \times n$ matrix onto an expanding $k$-dimensional Krylov subspace ($k \ll n$), the Lanczos algorithm rapidly computes extremal eigenvalues (the largest and smallest frequencies in the spectrum) using only one sparse matrix-vector product per step and minimal vector memory.
}

\subsubsection{The Three-Term Lanczos Recurrence}

\addThm{The Lanczos Factorization}{%
Let $\bA \in \fR^{n \times n}$ be real symmetric, and let $\bv_1 \in \fR^n$ be a normalized initial vector ($\|\bv_1\|_2 = 1$). Set $\bv_0 = \bzero$ and $\beta_0 = 0$.
For $j = 1, 2, \dots, k$:
\begin{enumerate}
    \item Compute the matrix-vector product: $\bw = \bA\bv_j$.
    \item Compute the diagonal Rayleigh quotient: $\alpha_j = \bv_j^T \bw$.
    \item Orthogonalize against the previous two vectors: $\bw \leftarrow \bw - \alpha_j \bv_j - \beta_{j-1} \bv_{j-1}$.
    \item Compute the subdiagonal norm: $\beta_j = \|\bw\|_2$.
    \item If $\beta_j = 0$, terminate. Otherwise, set $\bv_{j+1} = \bw / \beta_j$.
\end{enumerate}
In matrix notation, after $k$ steps, the orthonormal Lanczos basis $\bV_k = [\bv_1, \dots, \bv_k]$ satisfies the algebraic relation
\begin{align*}
    \bA \bV_k = \bV_k \bT_k + \beta_k \bv_{k+1} \be_k^T,
\end{align*}
where $\bT_k \in \fR^{k \times k}$ is the symmetric tridiagonal matrix
\begin{align*}
    \bT_k = \begin{pmatrix}
        \alpha_1 & \beta_1 & 0 & \dots & 0 \\
        \beta_1 & \alpha_2 & \beta_2 & \dots & 0 \\
        0 & \beta_2 & \alpha_3 & \ddots & \vdots \\
        \vdots & \vdots & \ddots & \ddots & \beta_{k-1} \\
        0 & 0 & \dots & \beta_{k-1} & \alpha_k
    \end{pmatrix}.
\end{align*}
}
\label{thm:lanczos-relation}

\addProp{Galerkin Projection Property}{%
In exact arithmetic, the Lanczos basis vectors are mutually orthonormal ($\bV_k^T\bV_k = \bI_k$). Multiplying the Lanczos relation on the left by $\bV_k^T$ yields the Rayleigh-Ritz Galerkin projection:
\begin{align*}
    \bT_k = \bV_k^T \bA \bV_k.
\end{align*}
}
\label{prop:lanczos-galerkin}

\subsubsection{Ritz Values, Ritz Vectors, and Error Bounds}

\addDef{Ritz Pairs}{%
Let $\bT_k \bs_i = \theta_i \bs_i$ be the eigendecomposition of the small symmetric tridiagonal matrix $\bT_k$, with normalized eigenvectors $\|\bs_i\|_2 = 1$.
\begin{enumerate}
    \item The eigenvalue $\theta_i \in \fR$ is termed the \textbf{$i$-th Ritz value} of $\bA$ with respect to $\mathcal{K}_k(\bA, \bv_1)$.
    \item The vector $\tilde{\bu}_i = \bV_k \bs_i \in \fR^n$ is termed the \textbf{$i$-th Ritz vector}.
\end{enumerate}
}
\label{def:ritz-pair}

\addThm{A Posteriori Residual Bound for Ritz Pairs}{%
The residual of the Ritz pair $(\theta_i, \tilde{\bu}_i)$ satisfies the explicit identity
\begin{align*}
    \|\bA \tilde{\bu}_i - \theta_i \tilde{\bu}_i\|_2 = \beta_k |s_{k,i}|,
\end{align*}
where $s_{k,i}$ is the final component of the $i$-th tridiagonal eigenvector $\bs_i \in \fR^k$.
}
\label{thm:ritz-residual}

\addPrf{%
Multiplying the Lanczos relation $\bA \bV_k = \bV_k \bT_k + \beta_k \bv_{k+1} \be_k^T$ on the right by $\bs_i$ gives
\begin{align*}
    \bA (\bV_k \bs_i) = \bV_k (\bT_k \bs_i) + \beta_k \bv_{k+1} (\be_k^T \bs_i)
    \implies \bA \tilde{\bu}_i = \theta_i \tilde{\bu}_i + \beta_k s_{k,i} \bv_{k+1}.
\end{align*}
Rearranging and taking the Euclidean norm yields
\begin{align*}
    \|\bA \tilde{\bu}_i - \theta_i \tilde{\bu}_i\|_2 = \|\beta_k s_{k,i} \bv_{k+1}\|_2 = \beta_k |s_{k,i}| \|\bv_{k+1}\|_2 = \beta_k |s_{k,i}|,
\end{align*}
because $\|\bv_{k+1}\|_2 = 1$. This allows monitoring the convergence of every Ritz pair for $O(k)$ operations without forming the $n$-dimensional Ritz vectors $\tilde{\bu}_i$.
}

\subsubsection{Loss of Orthogonality and Ghost Eigenvalues}

\addRemark{%
\textbf{(Paige's theorem on loss of orthogonality)}\enspace
In floating-point arithmetic, the three-term recurrence suffers from a fundamental numerical instability discovered by Chris Paige. The basis vectors $\bv_j$ do not remain orthogonal; rather, orthogonality is lost rapidly precisely as individual Ritz values converge to extreme eigenvalues of $\bA$. As a consequence, once an eigenvalue has converged, the algorithm continues to re-discover it in subsequent iterations, generating multiple spurious copies known as \textbf{ghost eigenvalues}.
}

\addRemark{%
\textbf{(Practical remedies for ghost eigenvalues)}\enspace
To maintain numerical reliability on large problems, scientific codes employ one of three strategies:
\begin{enumerate}
    \item \textbf{Full Reorthogonalization:} Reorthogonalize $\bv_{j+1}$ against all previous basis vectors $\bv_1, \dots, \bv_j$ using Modified Gram-Schmidt at every step ($O(k^2 n)$ cost).
    \item \textbf{Selective Reorthogonalization:} Orthogonalize only against converged Ritz vectors whenever loss of orthogonality is detected ($O(k \cdot n)$ cost).
    \item \textbf{Cullum-Willoughby Identification:} Run the unorthogonalized Lanczos iteration and filter out ghost eigenvalues by comparing the spectrum of $\bT_k$ with the spectrum of the submatrix formed by deleting the first row and column of $\bT_k$.
\end{enumerate}
}

\addExcr{%
\begin{enumerate}
    \item Execute two steps of the Lanczos algorithm by hand for the diagonal matrix $\bA = \operatorname{diag}(4, 2, 1)$ starting with the normalized vector $\bv_1 = \frac{1}{\sqrt{3}}(1, 1, 1)^T$. Explicitly write down $\alpha_1, \beta_1, \alpha_2$, and the $2 \times 2$ tridiagonal matrix $\bT_2$.
    \item For the $2 \times 2$ matrix $\bT_2$ computed above, find its Ritz values and evaluate the residual error bound using Thm~\ref{thm:ritz-residual}.
    \item Prove that the Lanczos algorithm applied to the shifted system $(\bA - \sigma\bI)$ produces the identical basis vectors $\bV_k$ as the original system, with tridiagonal matrix $\bT_k - \sigma\bI$.
\end{enumerate}
}
\label{ex:lanczos-exercises}

\end{document}
